\section{Fungsi Agregat: COUNT, SUM, AVG, MIN, MAX}

\textbf{Fungsi agregat} mengolah sekumpulan baris dan mengembalikan satu nilai ringkasan. Fungsi ini sangat berguna untuk analisis data: menghitung jumlah record, total penjualan, rata-rata harga, nilai minimum, dan maksimum. Fungsi agregat biasanya digunakan bersama \code{GROUP BY} untuk menghasilkan ringkasan per kelompok, namun juga dapat digunakan sendiri untuk ringkasan seluruh tabel \cite{mysqlDocs}.

\begin{table}[ht]
\centering
\begin{tabular}{|P{2.5cm}|P{5cm}|P{5cm}|}
\hline
\textbf{Fungsi} & \textbf{Deskripsi} & \textbf{Contoh} \\
\hline
\code{COUNT()} & Jumlah baris & \code{COUNT(*)} atau \code{COUNT(kolom)} \\
\hline
\code{SUM()} & Total nilai & \code{SUM(harga * stok)} \\
\hline
\code{AVG()} & Rata-rata nilai & \code{AVG(harga)} \\
\hline
\code{MIN()} & Nilai terkecil & \code{MIN(harga)} \\
\hline
\code{MAX()} & Nilai terbesar & \code{MAX(harga)} \\
\hline
\end{tabular}
\caption{Fungsi agregat utama di MySQL}
\end{table}

\begin{webcode}[language=SQL, caption={Penggunaan fungsi agregat}]
-- Jumlah total produk
SELECT COUNT(*) AS total_produk FROM produk;

-- Total nilai inventaris
SELECT SUM(harga * stok) AS total_inventaris FROM produk;

-- Rata-rata harga produk elektronik
SELECT AVG(harga) AS rata_harga
  FROM produk WHERE kategori = 'elektronik';

-- Produk termurah dan termahal
SELECT MIN(harga) AS termurah, MAX(harga) AS termahal FROM produk;

-- Ringkasan per kategori
SELECT kategori,
       COUNT(*) AS jumlah,
       AVG(harga) AS rata_harga,
       SUM(stok) AS total_stok
FROM produk GROUP BY kategori;
\end{webcode}

